\chap{3}{SYSTEM REQUIREMENTS SPECIFICATION}{3pt}
\vspace{-8pt}
\sect{3.1}{FUNCTIONAL REQUIREMENTS}
\vspace{2pt}
The functional requirements identify what the Evo2Serve system must do in order to operate as an AI-driven genetic variant analysis platform. They reflect real usage scenarios such as account access, genome browsing, gene and variant lookup, mutation entry, GPU-based scoring, comparison with ClinVar, calibration and verification.

Several of these requirements follow from the nature of genome data. Every position handled by the system must travel with the name of its assembly (GRCh37 or GRCh38) and its chromosome, because a coordinate is meaningless without them. Positions are entered and displayed one-based, as in ClinVar and in the HGVS nomenclature, and are converted to the zero-based intervals that the UCSC sequence service expects. A sequence consists of the bases A, C, G and T, and the character N where the base is unknown; a window that contains an N cannot be scored reliably and must be detected. Finally, the model needs a fixed amount of context around the variant, which the system supplies as a window of 8,193 bases centred on it.

\vspace{6pt}
\tsub{1. User Registration and Authentication}
\begin{subitemize}
\item The system shall allow a user to create an account with an e-mail address and a password, and to sign in and sign out.
\item The system shall issue a signed session token when a user signs in successfully, and shall reject a wrong password or an already registered e-mail address with a message.
\item Browsing shall be available without an account, whereas running the model shall require the user to be signed in.
\end{subitemize}

\vspace{4pt}
\tsub{2. Genome Browsing and Gene Search}
\begin{subitemize}
\item The system shall list the human genome assemblies published by UCSC (hg38 by default) and the chromosomes of the selected assembly.
\item Users shall be able to search genes by name (for example BRCA1) or browse the genes of a chromosome, and to open the gene details and the reference sequence.
\end{subitemize}

\vspace{4pt}
\tsub{3. Known Variant Retrieval from ClinVar}
\begin{subitemize}
\item The system shall retrieve the known variants of the selected gene from ClinVar using the position field of the selected assembly.
\item Each variant shall be shown with its type, position and clinical classification, and a single click shall run the model on it.
\end{subitemize}

\vspace{4pt}
\tsub{4. Variant Entry and Validation}
\begin{subitemize}
\item The user shall be able to enter a substitution, deletion, insertion or deletion--insertion by clicking a base or typing a position.
\item The system shall validate every request before scoring and reject an invalid request with a message; deletion and insertion lengths shall be limited to 500 bases.
\end{subitemize}

\tsub{5. Evo 2 Scoring Service}
\begin{subitemize}
\item The service shall fetch an 8\,kb context window from UCSC, build the mutated sequence, and score both sequences with Evo~2-7B.
\item It shall return the score difference $\Delta$, the label (``likely pathogenic'' or ``likely benign''), the confidence, the threshold used and a provisional flag.
\end{subitemize}

\vspace{4pt}
\tsub{6. Comparison with ClinVar}
\begin{subitemize}
\item The interface shall show the Evo~2 prediction and its confidence bar next to the ClinVar classification of the same variant.
\item It shall state whether the two agree and mark indel predictions as provisional.
\item While a prediction is running the interface shall show a loading state, and when a lookup or the model service fails it shall show an error message.
\end{subitemize}

\vspace{4pt}
\tsub{7. Calibration Routine}
\begin{subitemize}
\item An offline GPU job shall score the first 500 \textit{BRCA1} variants of the saturation workbook, compute the ROC curve of $-\Delta$ and the AUROC.
\item The job shall derive the Youden-optimal threshold and the class spreads used for the confidence value.
\end{subitemize}

\vspace{4pt}
\tsub{8. Verification and Reporting}
\begin{subitemize}
\item Scripts shall recompute the calibration data composition, the reference predictors, the assembly offset and the validator conformance.
\item Each script shall write its results to a machine-readable output file.
\end{subitemize}

\sectend{3.1}

\sect{3.2}{NON- FUNCTIONAL REQUIREMENTS}
Non-functional requirements are those that define the qualities of the system that are necessary to make it resilient, performant, and satisfactory for the users.

\tsub{1. Performance Requirements}
\begin{subitemize}
\item Web requests other than model inference shall complete within a second under typical loads.
\item A prediction requires two forward passes over about 8,200 bases; the interface shall stay responsive while it runs.
\end{subitemize}

\tsub{2. Reliability}
\begin{subitemize}
\item The GPU function shall retry failed calls twice, and a failure of UCSC, NCBI or the model service shall not crash the interface.
\end{subitemize}

\tsub{3. Scalability}
\begin{subitemize}
\item The GPU tier shall scale automatically from zero to three containers and scale down after 120 seconds of inactivity.
\item Model weights shall be cached on a persistent volume so that new containers do not download them again.
\end{subitemize}

\vspace{4pt}
\tsub{4. Usability and Interface Quality}
\begin{subitemize}
\item The interface shall be responsive, with a clear layout for assembly, gene, sequence, variants and results.
\item The interface shall use consistent labels and colour coding, for example for the ClinVar classes pathogenic, likely benign and uncertain significance.
\end{subitemize}

\tsub{5. Security and Data Integrity}
\begin{subitemize}
\item All external communication shall use HTTPS; secrets such as the signing key shall be held in environment variables and Modal secrets.
\item Passwords shall be stored only as hashes, the session token shall be held in an HTTP-only cookie, and every request shall be validated independently in the web tier and again in the GPU tier.
\item The GPU endpoint shall reject requests without a valid bearer token with HTTP 401.
\end{subitemize}

\tsub{6. Validity of Predictions}
\begin{subitemize}
\item The confidence value shall be documented as a distance heuristic and not as a probability.
\item Predictions shall be presented as computational research evidence and not as a clinical diagnosis.
\end{subitemize}

\tsub{7. Maintainability and Reproducibility}
\begin{subitemize}
\item The front end shall be organised in typed components and the back end in separate functions for windowing, editing and scoring.
\item The request schema shall be defined once per tier and tested against a shared case list.
\item Every number reported shall be reproducible without a GPU from the committed scripts and their machine-readable outputs.
\end{subitemize}

\sectend{3.2}

\sect{3.3}{HARDWARE REQUIREMENTS}
\vspace{2pt}
The system runs on three kinds of hardware: the developer machine, the cloud GPU that hosts the model, and the user's device. The specifications below are recommended for stable execution.

\vspace{4pt}
\sub{Minimum Requirements (For Development \& Testing):}
\begin{itemize}
\item \textbf{Processor:} Intel Core i3 or AMD equivalent (2.0 GHz or above)
\item \textbf{RAM:} 8 GB
\item \textbf{Storage:} 20 GB free SSD or HDD space
\item \textbf{Display:} 1366 x 768 resolution monitor or higher
\item \textbf{Network:} Reliable broadband connectivity
\end{itemize}
\sectend{3.3a}

\sub{Recommended Requirements (For Optimal Performance \& Deployment):}
\begin{itemize}
\item \textbf{Processor:} Intel Core i5/i7 or AMD Ryzen 5/7 (3.0 GHz or above)
\item \textbf{RAM:} 16 GB or higher
\item \textbf{Storage:} 100 GB SSD or higher
\item \textbf{Display:} Full HD (1920 x 1080)
\item \textbf{Network:} Speedy broadband (UCSC, NCBI and GPU endpoint calls)
\item \textbf{GPU Server (Cloud):} NVIDIA H100 GPU provided by Modal for Evo~2-7B inference
\end{itemize}

\vspace{4pt}
\sub{Optional / Advanced Hardware}
\begin{itemize}
\item Additional H100-class GPUs for the Evo~2-40B model, which requires several GPUs
\item Persistent cloud storage volume for the Hugging Face model cache
\item Redundant storage for backups of the user database
\end{itemize}
\sectend{3.3}

\sect{3.4}{SOFTWARE REQUIREMENTS}
The software environment plays a crucial role in the design, development, and execution of the genetic variant analysis system. The following software components are required:

\vspace{4pt}
\sub{Operating System:}
\begin{itemize}
\item Windows 10/11
\item Linux (Ubuntu 22.04 or later; the GPU container image is based on Ubuntu 22.04)
\item macOS (for cross-platform development)
\end{itemize}

\vspace{4pt}
\sub{Programming Languages \& Frameworks:}
\begin{itemize}
\item \textbf{Python 3.12} -- GPU service, calibration job and analysis scripts
\item \textbf{FastAPI and Pydantic} -- REST endpoint and request validation
\item \textbf{TypeScript and Node.js 22} -- front end, API routes and the Zod validation script
\item \textbf{Next.js 15 and React 19} -- web application framework (T3 stack)
\end{itemize}
\sectend{3.4a}

\sub{Machine Learning \& Data Science Libraries:}
\begin{itemize}
\item \textbf{Evo 2 package:} Official implementation of the model; loads the \texttt{evo2\_7b} checkpoint and provides \texttt{score\_sequences}.
\item \textbf{PyTorch 2.7.1 and FlashAttention 2.8.0:} Deep learning runtime built for CUDA 12.6.
\item \textbf{pandas, scikit-learn and Biopython:} Workbook handling, ROC analysis and FASTA parsing.
\item \textbf{matplotlib and seaborn:} Plots produced by the calibration job.
\end{itemize}

\vspace{4pt}
\sub{Frontend \& Interface Components:}
\begin{itemize}
\item \textbf{Tailwind CSS 4 and shadcn/ui:} Styling and accessible interface components.
\item \textbf{Zod:} Schema validation of the analysis request in the web tier.
\item \textbf{jose and bcryptjs:} Signed session tokens and password hashing.
\item \textbf{Mongoose and MongoDB:} User account storage.
\end{itemize}

\vspace{4pt}
\sub{Platforms, Data \& External APIs:}
\begin{itemize}
\item \textbf{Modal:} Serverless GPU platform that hosts the FastAPI endpoint and the calibration job.
\item \textbf{UCSC Genome Browser REST API:} Assemblies, chromosomes and reference sequence.
\item \textbf{NCBI Gene, E-utilities and ClinVar:} Gene search, gene details and known variants.
\item \textbf{BRCA1 saturation genome editing workbook:} Experimental labels for 3,893 variants.
\end{itemize}

\vspace{4pt}
\sub{Developer Tools:}
\begin{itemize}
\item \textbf{Visual Studio Code:} Primary development environment.
\item \textbf{Git \& GitHub:} Version control and repository hosting.
\item \textbf{Modal CLI and npm:} Deployment of the GPU service and management of web dependencies.
\item \textbf{Excalidraw and \LaTeX:} Architecture sketches and report preparation.
\end{itemize}
\sectend{3.4}

\vspace{6pt}
\sect{3.5}{SYNOPSIS}
\vspace{4pt}
\textbf{Evo2Serve} is an AI-driven genetic variant analysis platform that integrates genome browsing, ClinVar lookup, GPU-based zero-shot scoring with Evo~2-7B, comparison with expert classifications, authenticated access and a verification suite within a unified system. The platform supports four mutation types, an 8\,kb context window, a calibrated decision threshold and a flagged confidence value. Its non-functional requirements prioritise response behaviour, reliability, architectural scalability, security, interface usability, validity of predictions and maintainability.

The functional requirements outline the capabilities of the system: account access, assembly and gene browsing, retrieval of known variants, validated variant entry, scoring on a serverless GPU, a comparison view, an offline calibration job and reproducible verification scripts. These modules ensure that users can obtain a model-based opinion on a mutation without any installation and can see it next to curated evidence.

The non-functional requirements focus on a responsive interface, retries and clear error reporting, automatic scaling between zero and three GPU containers, token-based security on both tiers, honest wording of predictions and a modular, typed code base.

Built on a software stack of Python, FastAPI, Next.js, React, PyTorch and Modal, with public genome and variant services as data sources, \textbf{Evo2Serve} is designed to run on a standard development machine while the model runs on a cloud H100 GPU. By establishing clear functional boundaries, quality attributes and operating environments, this requirements specification provides a foundation for the system architecture, the implementation and the verification described in the following chapters.
\sectend{3.5}
